\documentclass[11pt,a4paper]{article}

% ==========  Packages  ==========
\usepackage[margin=1in]{geometry}
\usepackage{setspace}
\usepackage{parskip}
\usepackage{xcolor}
\usepackage{listings}
\usepackage{booktabs}
\usepackage{array}
\usepackage{enumitem}
\usepackage{tcolorbox}
\usepackage[colorlinks=true,linkcolor=blue!60!black,urlcolor=blue!60!black,citecolor=blue!60!black]{hyperref}

\setstretch{1.15}

% ==========  Colors  ==========
\definecolor{pwshBg}{HTML}{F4F6FA}
\definecolor{pwshKw}{HTML}{0050A0}
\definecolor{pwshStr}{HTML}{8B2C00}
\definecolor{pwshCom}{HTML}{4F7F4F}
\definecolor{pwshVar}{HTML}{7B2D8E}

\definecolor{bashBg}{HTML}{F8F1E4}
\definecolor{bashKw}{HTML}{8B2C00}
\definecolor{bashStr}{HTML}{4F7F4F}
\definecolor{bashCom}{HTML}{888888}

% ==========  PowerShell language for listings  ==========
\lstdefinelanguage{PowerShell}{
  morekeywords={if,else,elseif,foreach,for,while,do,switch,return,function,filter,break,continue,
    param,try,catch,finally,throw,begin,process,end,in,not,and,or,
    Get-ChildItem,Set-Location,Get-Location,Get-Content,Set-Content,Add-Content,Out-File,
    Where-Object,Select-Object,ForEach-Object,Sort-Object,Group-Object,Measure-Object,
    Import-Csv,Export-Csv,ConvertFrom-Csv,ConvertTo-Csv,ConvertFrom-Json,ConvertTo-Json,
    Invoke-RestMethod,Invoke-WebRequest,Get-Command,Get-Help,Get-Member,Get-Alias,
    New-Item,Remove-Item,Copy-Item,Move-Item,Test-Path,Rename-Item,Resolve-Path,
    Get-Process,Start-Process,Stop-Process,Write-Host,Write-Output,Write-Error,Write-Warning,
    Read-Host,Out-Null,Out-String,Out-GridView,Format-Table,Format-List,
    Set-ExecutionPolicy,Get-ExecutionPolicy,Get-Item,Get-Date,Get-Random,
    New-Object,Select-String,Start-Sleep,Tee-Object},
  sensitive=false,
  morecomment=[l]{\#},
  morestring=[b]",
  morestring=[b]'
}

% ==========  Listing styles  ==========
\lstdefinestyle{pwshstyle}{
  language=PowerShell,
  backgroundcolor=\color{pwshBg},
  basicstyle=\ttfamily\small,
  keywordstyle=\color{pwshKw}\bfseries,
  stringstyle=\color{pwshStr},
  commentstyle=\color{pwshCom}\itshape,
  showstringspaces=false,
  breaklines=true,
  breakatwhitespace=true,
  frame=single,
  framesep=4pt,
  rulecolor=\color{pwshKw!30},
  xleftmargin=8pt,
  xrightmargin=4pt,
  aboveskip=8pt,
  belowskip=8pt,
  columns=fullflexible,
  keepspaces=true
}

\lstdefinestyle{bashstyle}{
  language=bash,
  backgroundcolor=\color{bashBg},
  basicstyle=\ttfamily\small,
  keywordstyle=\color{bashKw}\bfseries,
  stringstyle=\color{bashStr},
  commentstyle=\color{bashCom}\itshape,
  showstringspaces=false,
  breaklines=true,
  breakatwhitespace=true,
  frame=single,
  framesep=4pt,
  rulecolor=\color{bashKw!30},
  xleftmargin=8pt,
  xrightmargin=4pt,
  aboveskip=8pt,
  belowskip=8pt,
  columns=fullflexible,
  keepspaces=true
}

\lstset{style=pwshstyle}

% ==========  Macros  ==========
\newcommand{\cmd}[1]{\texttt{\textcolor{pwshKw}{#1}}}
\newcommand{\bcmd}[1]{\texttt{\textcolor{bashKw}{#1}}}

% A "compare with what you know" callout box
\newtcolorbox{compare}[1][]{
  colback=yellow!8,
  colframe=yellow!50!orange,
  fonttitle=\bfseries,
  title=Compare with what you know,
  sharp corners,
  boxrule=0.4pt,
  left=8pt,right=8pt,top=4pt,bottom=4pt,
  #1
}

\newtcolorbox{gotcha}[1][]{
  colback=red!5,
  colframe=red!50,
  fonttitle=\bfseries,
  title=Gotcha,
  sharp corners,
  boxrule=0.4pt,
  left=8pt,right=8pt,top=4pt,bottom=4pt,
  #1
}

% ==========  Title  ==========
\title{\bfseries PowerShell for the Data \& ML Engineer\\
\large A practical tutorial for Python / cron / AWS / Mongo / MySQL users\\moving onto Windows}
\author{}
\date{}

\begin{document}
\maketitle

\begin{abstract}
\noindent
This tutorial teaches the PowerShell you actually need as a data engineer or
machine-learning engineer on Windows: navigating the filesystem, transforming
data with the object pipeline, parsing CSV and JSON, calling REST APIs,
managing environment variables for AWS / MongoDB / MySQL credentials, and
running Python with \texttt{venv} and \texttt{uv}. Each concept is introduced
through a worked example, with a side-by-side comparison to the bash or
Python equivalent you already know. PowerShell 7+ is assumed throughout
(the cross-platform \texttt{pwsh}); 5.1-only quirks are flagged when they
matter.
\end{abstract}

\tableofcontents
\clearpage

% =========================================================
\section{What PowerShell Is (and Isn't)}
% =========================================================

The one idea to grasp before you write any PowerShell:
\textbf{PowerShell pipes objects, not text.}

In bash, every command emits a stream of bytes, and the next command in the
pipeline parses those bytes back into a structure. If you want the size of
every file over 1\,MB, you run \bcmd{ls -l}, then \bcmd{awk} to chop columns,
then \bcmd{sort -n}, and a typo in the column index silently breaks
everything. In PowerShell, \cmd{Get-ChildItem} returns real
\texttt{FileInfo} objects with properties like \texttt{Length}, \texttt{Name},
\texttt{LastWriteTime}, and the next stage of the pipeline filters and
sorts by those properties directly. There is no parsing step, and no
formatting choices to break.

\begin{lstlisting}[style=pwshstyle]
# All files over 1 MB in the current directory, biggest first.
Get-ChildItem |
    Where-Object Length -gt 1MB |
    Sort-Object Length -Descending |
    Select-Object Name, Length
\end{lstlisting}

\begin{lstlisting}[style=bashstyle]
# Roughly the same in bash -- and notice how much
# more fragile the column parsing is.
ls -la | awk '$5 > 1048576 {print $9, $5}' | sort -k2 -nr
\end{lstlisting}

\subsection*{The Verb-Noun convention}

Every cmdlet (PowerShell's word for a command) is named
\texttt{Verb-Noun}: \cmd{Get-ChildItem}, \cmd{Set-Location},
\cmd{New-Item}, \cmd{Remove-Item}, \cmd{Import-Csv},
\cmd{Invoke-RestMethod}. The verbs come from a small, deliberate
vocabulary (\texttt{Get}, \texttt{Set}, \texttt{New}, \texttt{Remove},
\texttt{Add}, \texttt{Test}, \texttt{Invoke}, \texttt{Out},
\texttt{ConvertTo}, \texttt{ConvertFrom}, \dots). Once you know roughly
what you want to do, you can often guess the cmdlet name.

\subsection*{Cmdlets, aliases, and discovery}

Most Unix command names work in PowerShell as \emph{aliases} for
real cmdlets. \texttt{ls} is an alias for \cmd{Get-ChildItem};
\texttt{cd} is an alias for \cmd{Set-Location}; \texttt{cat} is
\cmd{Get-Content}; \texttt{rm} is \cmd{Remove-Item}. Aliases are
fine in the interactive shell, but \textbf{use the full cmdlet name
in scripts} so the script is portable and self-documenting.

Three cmdlets you will use constantly to learn the rest:

\begin{lstlisting}[style=pwshstyle]
Get-Command *csv*           # find every cmdlet whose name contains "csv"
Get-Help Import-Csv -Full   # the full help page (parameters, examples)
Get-Item .\data.csv | Get-Member   # list every property and method of an object
\end{lstlisting}

\cmd{Get-Member} is the one to reach for first: \emph{pipe anything to
it} and PowerShell lists the properties and methods that object
exposes. This is how you find out, for example, that the result of
\cmd{Invoke-RestMethod} on a JSON endpoint is already a real object with
dot-accessible properties.

\begin{compare}
In Python, you call \texttt{dir(obj)} or \texttt{help(obj)} to explore an
object's interface. \cmd{Get-Member} is the same idea, with the bonus
that it works inline at the end of a pipeline.
\end{compare}

% =========================================================
\section{Getting Around: Files and Directories}
% =========================================================

The cmdlets in this section replace the daily-driver Unix commands.
They all accept paths with spaces if you quote them, and they all
take \texttt{-Recurse} and \texttt{-Force} flags where you would
expect.

\begin{lstlisting}[style=pwshstyle]
Get-Location                                # pwd
Set-Location C:\Users\shiru\Desktop         # cd
Get-ChildItem                                # ls
Get-ChildItem -Recurse -Filter *.parquet    # find . -name "*.parquet"
Get-ChildItem -Force                         # ls -a (include hidden files)

New-Item -ItemType Directory data\raw       # mkdir data\raw
New-Item -ItemType Directory -Force a\b\c   # mkdir -p a\b\c
New-Item -ItemType File notes.txt           # touch notes.txt

Copy-Item .\data.csv .\backup\data.csv
Copy-Item -Recurse .\src .\src_backup       # cp -r
Move-Item .\old_name.csv .\new_name.csv     # mv (also rename)
Remove-Item .\tmp.csv                       # rm
Remove-Item -Recurse -Force .\.venv         # rm -rf

Test-Path .\data.csv                         # [ -e data.csv ]; returns $true/$false
Resolve-Path .\data.csv                      # absolute path
\end{lstlisting}

\subsection*{Paths and quoting}

PowerShell paths use backslash on Windows (\texttt{C:\textbackslash{}Users\textbackslash{}shiru\textbackslash{}data.csv})
but forward slash works too. Surround paths containing spaces with
double quotes:

\begin{lstlisting}[style=pwshstyle]
Get-Content "C:\Users\shiru\My Documents\notes.txt"
\end{lstlisting}

The backtick (\texttt{`}) is the escape character (\emph{not} backslash),
and the dollar sign (\texttt{\$}) is special inside double quotes (it
interpolates). Inside single quotes, nothing is special:

\begin{lstlisting}[style=pwshstyle]
Write-Output "Today is $env:USERNAME"     # interpolates -> "Today is shiru"
Write-Output 'Today is $env:USERNAME'     # literal      -> "Today is $env:USERNAME"
Write-Output "Price is `$5"               # backtick-escape -> "Price is $5"
\end{lstlisting}

\subsection*{Safety: -WhatIf and -Confirm}

Every destructive cmdlet (\cmd{Remove-Item}, \cmd{Stop-Process},
\cmd{Set-Content} on an existing file, \dots) accepts
\texttt{-WhatIf} to \emph{print what would happen without doing it},
and \texttt{-Confirm} to prompt before each action. Get in the habit
of running anything you're not sure about with \texttt{-WhatIf} first:

\begin{lstlisting}[style=pwshstyle]
Remove-Item -Recurse -Force .\.venv -WhatIf
# What if: Performing the operation "Remove Directory" on target ".venv".
\end{lstlisting}

\begin{compare}
Unix has nothing equivalent. The closest is running a \texttt{find}
command before a \texttt{find -delete} to inspect the matches. With
\texttt{-WhatIf} every destructive cmdlet has dry-run built in.
\end{compare}

% =========================================================
\section{The Object Pipeline}
% =========================================================

This is the most useful section in the tutorial. If you only learn
five PowerShell cmdlets, make them these:

\begin{center}
\begin{tabular}{ll}
\toprule
Cmdlet & What it does \\
\midrule
\cmd{Where-Object}    & filter (\texttt{WHERE}, \texttt{df.query}, \texttt{filter()})  \\
\cmd{Select-Object}   & project columns / take N (\texttt{SELECT}, \texttt{df[[...]]}) \\
\cmd{ForEach-Object}  & map / transform each item (\texttt{map}, \texttt{apply}) \\
\cmd{Sort-Object}     & sort by a property (\texttt{ORDER BY}, \texttt{df.sort\_values}) \\
\cmd{Group-Object}    & group by a key (\texttt{GROUP BY}, \texttt{df.groupby}) \\
\cmd{Measure-Object}  & count / sum / avg / min / max (\texttt{aggregate}) \\
\bottomrule
\end{tabular}
\end{center}

The objects flow left-to-right; the variable \texttt{\$\_} refers to
the current object inside a \texttt{\{ ... \}} script block.

\subsection*{Worked example: log file analysis}

Suppose you have a directory of \texttt{*.log} files and want to know
which file has the most lines mentioning \texttt{ERROR}.

\begin{lstlisting}[style=pwshstyle]
Get-ChildItem -Filter *.log |
    ForEach-Object {
        [PSCustomObject]@{
            File   = $_.Name
            Errors = (Select-String -Path $_.FullName -Pattern 'ERROR').Count
        }
    } |
    Sort-Object Errors -Descending |
    Select-Object -First 5
\end{lstlisting}

This is the entire pipeline, end to end: enumerate files, build a small
object per file with \texttt{File} and \texttt{Errors} properties, sort
by error count, take the top 5.

\subsection*{Block form vs.\ property form}

\cmd{Where-Object} has two syntaxes. They do the same thing; the
\emph{property form} is shorter and works for simple comparisons.

\begin{lstlisting}[style=pwshstyle]
# Property form (simple comparison).
Get-ChildItem | Where-Object Length -gt 1MB

# Block form (any expression).
Get-ChildItem | Where-Object { $_.Length -gt 1MB -and $_.Name -like '*.csv' }
\end{lstlisting}

Operators are spelled out (not symbolic) because \texttt{<} and \texttt{>}
are reserved for redirection. The full set you'll use:

\begin{center}
\renewcommand{\arraystretch}{1.1}
\begin{tabular}{p{0.18\linewidth} p{0.74\linewidth}}
\toprule
\textbf{Operator} & \textbf{Meaning} \\
\midrule
\texttt{-eq}, \texttt{-ne}            & equal, not equal \\
\texttt{-lt}, \texttt{-le}            & less than, less or equal \\
\texttt{-gt}, \texttt{-ge}            & greater than, greater or equal \\
\texttt{-like}                        & wildcard match (\texttt{*}, \texttt{?}) -- case insensitive \\
\texttt{-clike}                       & wildcard, case sensitive (prefix \texttt{c} on any op makes it case-sensitive) \\
\texttt{-match}                       & regex match (populates \texttt{\$Matches}) \\
\texttt{-notmatch}                    & negated regex \\
\texttt{-in}, \texttt{-notin}          & membership: \texttt{\$x -in @(1,2,3)} \\
\texttt{-contains}, \texttt{-notcontains} & reverse membership: collection \texttt{-contains \$x} \\
\texttt{-and}, \texttt{-or}, \texttt{-not}, \texttt{-xor} & boolean combinators \\
\texttt{-is}, \texttt{-isnot}          & type check: \texttt{\$x -is [string]} \\
\bottomrule
\end{tabular}
\end{center}

These operators broadcast naturally over collections, so
\texttt{@(1,2,3,4,5) -gt 2} returns \texttt{@(3,4,5)} -- the same
trick lets you do quick set filters without an explicit
\cmd{Where-Object} call.

\subsection*{Picking columns and rows with \texttt{Select-Object}}

\cmd{Select-Object} chooses which properties to keep and how many
objects to pass on. A few forms you'll use constantly:

\begin{lstlisting}[style=pwshstyle]
Get-Process | Select-Object Name, CPU, Id      # keep only these properties
Get-Process | Select-Object -First 5           # first 5 objects (like head)
Get-Process | Select-Object -Last 5            # last 5 (like tail)
Get-ChildItem | Select-Object -Unique Extension # drop duplicate values

# -ExpandProperty unwraps a single property to its raw value (not an object):
Get-Process | Select-Object -ExpandProperty Name   # plain strings, not objects

# A calculated property adds a column computed on the fly:
Get-ChildItem -File |
    Select-Object Name,
                  @{ Name = 'SizeKB'; Expression = { [math]::Round($_.Length / 1KB, 1) } }
\end{lstlisting}

The \texttt{@\{ Name = ...; Expression = \{ ... \} \}} hashtable is how
you rename or derive a column -- the rough equivalent of pandas'
\texttt{df.assign(SizeKB = df.Length / 1024)}.

\subsection*{Transforming each item with \texttt{ForEach-Object}}

\cmd{ForEach-Object} runs a script block once per item. In PowerShell 7+
you can also pass a property or method name directly, which is shorter:

\begin{lstlisting}[style=pwshstyle]
1..3 | ForEach-Object { $_ * $_ }          # 1, 4, 9
Get-ChildItem *.log | ForEach-Object Name  # shorthand: just the Name property
'a.txt', 'b.txt' | ForEach-Object { $_.ToUpper() }  # call a method per item

# -Begin runs once before, -End once after -- handy for running totals:
Get-ChildItem -File |
    ForEach-Object -Begin { $total = 0 } `
                   -Process { $total += $_.Length } `
                   -End { "Total bytes: $total" }
\end{lstlisting}

\subsection*{Group, count, summarize}

Counting events by category is one line:

\begin{lstlisting}[style=pwshstyle]
# How many files per extension in this directory tree?
Get-ChildItem -Recurse -File |
    Group-Object Extension |
    Sort-Object Count -Descending |
    Select-Object Count, Name

# -NoElement keeps just the counts (drops the grouped objects -- much lighter):
Get-ChildItem -Recurse -File | Group-Object Extension -NoElement
\end{lstlisting}

\cmd{Measure-Object} does numeric aggregation -- pass the flags for the
statistics you want:

\begin{lstlisting}[style=pwshstyle]
Get-ChildItem -Recurse -File |
    Measure-Object Length -Sum -Average -Maximum -Minimum
# Count   : 1842
# Average : 28734.21
# Sum     : 52928711
# Maximum : 14523904
# Minimum : 12

# It also counts lines, words, and characters in text:
Get-Content .\app.log | Measure-Object -Line -Word -Character
\end{lstlisting}

\subsection*{Parallel execution (PowerShell 7+)}

\cmd{ForEach-Object} accepts a \texttt{-Parallel} switch in
PowerShell 7+ that runs the script block on multiple threads -- ideal
for I/O-bound work like hitting hundreds of API endpoints or reading
many files at once.

\begin{lstlisting}[style=pwshstyle]
# Fetch metadata for 100 URLs in parallel, 10 at a time.
$urls = Get-Content .\urls.txt
$urls | ForEach-Object -Parallel {
    $r = Invoke-RestMethod -Uri $_
    [PSCustomObject]@{ url = $_; status = $r.status }
} -ThrottleLimit 10 |
    Export-Csv -NoTypeInformation -Path .\statuses.csv -Encoding utf8
\end{lstlisting}

Two things to know:

\begin{itemize}[leftmargin=*]
\item Inside \texttt{-Parallel}, the variable \texttt{\$\_} still
refers to the current item, but \emph{outside-scope variables are
not visible} -- pass them in with \texttt{\$using:varname}.
\item Order is not preserved. Use it for embarrassingly parallel
work; for sequential ETL where each step depends on the previous,
stick with plain \cmd{ForEach-Object}.
\end{itemize}

\begin{compare}
The PowerShell pipeline is so close to a pandas method chain that the
mapping is almost mechanical:
\[
\texttt{df[df.size > 1e6].sort\_values('size', ascending=False).head(5)}
\]
becomes
\[
\cmd{Get-ChildItem}\;|\;\cmd{Where-Object}\;\texttt{Length -gt 1MB}\;|\;\cmd{Sort-Object}\;\texttt{Length -Descending}\;|\;\cmd{Select-Object}\;\texttt{-First 5}.
\]
\end{compare}

% =========================================================
\section{Variables and Data Types}
% =========================================================

\subsection*{Variables}

Variable names start with \texttt{\$}. Assignment uses \texttt{=}; types
are inferred but you can cast.

\begin{lstlisting}[style=pwshstyle]
$name = 'shiru'
$count = 42
$ratio = 0.95
$today = Get-Date

# Casting:
$n = [int]"7"           # 7
$x = [double]"3.14"     # 3.14
$d = [datetime]"2026-01-15"
\end{lstlisting}

PowerShell variables are not scoped to a block by default (unlike
Python's function-scoped locals); they live in the enclosing scope.
For scripts, this is rarely an issue, but be aware of it when you
write loops.

\subsection*{Arrays}

The \texttt{@(...)} syntax builds an array; you can also build one
implicitly by listing items separated by commas. Indexing is
zero-based and uses square brackets, exactly like Python.

\begin{lstlisting}[style=pwshstyle]
$nums = @(1, 2, 3, 4, 5)
$nums[0]        # 1
$nums[-1]       # 5  (negative indexing works)
$nums[1..3]     # 2, 3, 4  (range slicing)
$nums.Length    # 5
$nums += 6      # append (creates a NEW array; arrays are immutable)

# Iterate:
foreach ($n in $nums) { Write-Output ($n * 2) }

# Or use the pipeline:
$nums | ForEach-Object { $_ * 2 }
\end{lstlisting}

\begin{gotcha}
The \texttt{+=} on arrays is \emph{not} cheap. It creates a brand-new
array each time. For large collections, build a
\texttt{[System.Collections.Generic.List[object]]} instead, or
use the pipeline to emit items.
\end{gotcha}

\subsection*{Hashtables (dicts)}

Hashtables are PowerShell's dictionary type.

\begin{lstlisting}[style=pwshstyle]
$user = @{
    name = 'shiru'
    role = 'Data Engineer'
    languages = @('Python', 'SQL', 'PowerShell')
}

$user.name              # 'shiru'
$user['role']           # 'Data Engineer' (both syntaxes work)
$user.email = 'shirui7352136@gmail.com'   # add a key
$user.ContainsKey('name')                  # $true
$user.Keys
$user.Values
\end{lstlisting}

\begin{compare}
Python dict literal \texttt{\{"name": "shiru", "role": "DE"\}} is
PowerShell's \texttt{@\{name='shiru'; role='DE'\}}. Note the
semicolons (or newlines) between entries -- not commas.
\end{compare}

\subsection*{Custom objects}

Cmdlets in a pipeline want objects, not bare hashtables. Promote a
hashtable to an object with \texttt{[PSCustomObject]}:

\begin{lstlisting}[style=pwshstyle]
$row = [PSCustomObject]@{
    user_id = 42
    score   = 0.91
    label   = 'churn'
}
$row | Get-Member        # see the properties
$row | Export-Csv -NoTypeInformation -Path .\out.csv
\end{lstlisting}

This is how you build CSV-shaped data on the fly and pipe it into the
data-wrangling cmdlets you'll meet in Section~\ref{sec:csv}.

\subsection*{Booleans and \$null}

\texttt{\$true}, \texttt{\$false}, and \texttt{\$null} are the three
constants you need. The comparison operators short-circuit (\texttt{-and},
\texttt{-or}, \texttt{-not}). Anything zero-length, zero-valued, or
\texttt{\$null} is falsy in conditions:

\begin{lstlisting}[style=pwshstyle]
if ($name) { Write-Output "got a name" }
if ($null -eq $row.email) { Write-Output "email is missing" }
\end{lstlisting}

% =========================================================
\section{Strings and Formatting}
% =========================================================

PowerShell strings come in four useful forms. Picking the right one
avoids a common Windows headache: escapes that get mangled.

\begin{lstlisting}[style=pwshstyle]
$name = 'shiru'

"Hello $name"              # double-quoted: variables interpolate
'Hello $name'              # single-quoted: literal, no interpolation
"Sum: $($a + $b)"          # subexpression form for complex inserts
"Path: $env:USERPROFILE"   # env: drive interpolates too

# Here-strings preserve newlines and ignore most quoting issues:
$body = @"
{
    "user": "$name",
    "count": 42
}
"@                          # closing "@ MUST be at column 0

# Single-quoted here-string -- literal, no interpolation, ideal for
# regex patterns or JSON templates with literal $ in them:
$pattern = @'
^\$[0-9]+\.[0-9]{2}$
'@
\end{lstlisting}

\begin{gotcha}
Inside a here-string, the closing \texttt{"@} or \texttt{'@} must be
at the start of the line, with no leading whitespace. Indenting it
is a parse error.
\end{gotcha}

\subsection*{Formatting numbers and dates}

The \texttt{-f} operator is the printf-style formatter. It uses .NET
format specifiers.

\begin{lstlisting}[style=pwshstyle]
"{0:N2}"   -f 1234.5678              # "1,234.57"  (2-decimal, with separators)
"{0:P1}"   -f 0.952                  # "95.2 %"     (percent)
"{0:F4}"   -f (1/3)                  # "0.3333"     (fixed-point)
"{0:yyyy-MM-dd}" -f (Get-Date)       # "2026-05-25"
"{0,-10} {1,5}" -f 'shiru', 42       # "shiru          42"  (padded columns)
\end{lstlisting}

\subsection*{Regex on strings}

\begin{lstlisting}[style=pwshstyle]
'AWS-prod-eu-west-1' -match '^AWS-(prod|stage)-([a-z0-9-]+)$'
# $true -- and the captures land in $Matches:
$Matches[1]   # 'prod'
$Matches[2]   # 'eu-west-1'

'foo,bar,baz' -split ','                # @('foo','bar','baz')
'foo bar baz' -replace '\s+', '_'       # 'foo_bar_baz'

# Multiline regex on a file:
Select-String -Path .\app.log -Pattern '^\d{4}-\d{2}-\d{2}.*ERROR' |
    Select-Object LineNumber, Line
\end{lstlisting}

\begin{compare}
\texttt{-replace} is regex by default. To do a literal replacement,
use the \texttt{.Replace()} string method:
\texttt{'a.b.c'.Replace('.', '/')}.
\end{compare}

% =========================================================
\section{File I/O -- Plain Text}
% =========================================================

\cmd{Get-Content} reads, \cmd{Set-Content} writes (overwrite),
\cmd{Add-Content} appends.

\begin{lstlisting}[style=pwshstyle]
# Read whole file into a single string vs. array of lines.
$lines = Get-Content .\app.log               # array, one element per line
$blob  = Get-Content .\app.log -Raw          # single string with embedded newlines

# Write / overwrite a file.
Set-Content -Path .\notes.txt -Value "First line"
Set-Content -Path .\notes.txt -Value @('First', 'Second', 'Third')   # array -> multi-line

# Append.
Add-Content -Path .\notes.txt -Value "Another line"

# Stream a large file line-by-line without loading it all into memory.
Get-Content .\huge.log -ReadCount 0 | ForEach-Object {
    if ($_ -match 'ERROR') { $_ }
}
\end{lstlisting}

\subsection*{Encoding}

A frequent source of cross-platform bugs: \cmd{Set-Content} and
\cmd{Out-File} write UTF-8 \emph{with BOM} by default in PowerShell 5.1
and UTF-8 \emph{without BOM} in PowerShell 7+. If you're writing data
that will be read by a non-Windows tool (Python's
\texttt{open(\dots, encoding='utf-8')}, AWS S3, jq), be explicit:

\begin{lstlisting}[style=pwshstyle]
Set-Content -Path .\out.txt -Value $text -Encoding utf8       # UTF-8 without BOM on PS7
Set-Content -Path .\out.txt -Value $text -Encoding utf8NoBOM  # explicit, PS7 only
Set-Content -Path .\out.txt -Value $text -Encoding ascii      # ASCII
\end{lstlisting}

\begin{gotcha}
\cmd{Out-File} has different default encoding from \cmd{Set-Content}
historically. In any code that crosses tool boundaries, \textbf{always
pass \texttt{-Encoding utf8}} explicitly.
\end{gotcha}

% =========================================================
\section{Data Wrangling -- CSV}
\label{sec:csv}
% =========================================================

This is where PowerShell starts to feel good for data work.
\cmd{Import-Csv} reads a CSV and gives back \textbf{one object per row},
with the column names as properties. No manual splitting, no quoting
bugs, no off-by-one column indexes.

Example file \texttt{users.csv}:
\begin{lstlisting}[style=pwshstyle,language=]
user_id,name,role,active,score
1,Alice,Data Engineer,true,0.91
2,Bob,ML Engineer,true,0.87
3,Carol,Analyst,false,0.62
4,Dan,Data Engineer,true,0.78
\end{lstlisting}

\subsection*{Read, filter, project, write}

\begin{lstlisting}[style=pwshstyle]
$users = Import-Csv .\users.csv

# Filter rows where active is "true" and score >= 0.8.
$top = $users |
    Where-Object { $_.active -eq 'true' -and [double]$_.score -ge 0.8 } |
    Select-Object name, role, score

$top | Format-Table          # pretty print
$top | Export-Csv -NoTypeInformation -Path .\top.csv -Encoding utf8

# Add a derived column with a calculated property, then write it out:
$users |
    Select-Object name, role,
                  @{ Name = 'score_pct'; Expression = { [int]([double]$_.score * 100) } } |
    Export-Csv -NoTypeInformation -Path .\users_pct.csv -Encoding utf8
\end{lstlisting}

Two things to notice:

\begin{enumerate}[leftmargin=*]
\item Every column comes back as a \emph{string}, even \texttt{score}.
You have to cast with \texttt{[double]} (or \texttt{[int]},
\texttt{[bool]}, \texttt{[datetime]}) before numeric comparisons.
\item \texttt{-NoTypeInformation} is critical when writing -- otherwise
\cmd{Export-Csv} prepends an annoying \texttt{\#TYPE} comment line
that confuses every non-PowerShell consumer.
\end{enumerate}

\subsection*{Group by, aggregate}

\begin{lstlisting}[style=pwshstyle]
# Average score per role.
Import-Csv .\users.csv |
    Group-Object role |
    ForEach-Object {
        [PSCustomObject]@{
            role         = $_.Name
            n            = $_.Count
            avg_score    = ($_.Group | Measure-Object {[double]$_.score} -Average).Average
        }
    } |
    Sort-Object avg_score -Descending |
    Format-Table
\end{lstlisting}

\subsection*{Joining two CSVs}

PowerShell has no SQL-style \texttt{JOIN}, but the hashtable-as-index
trick gets you the same answer with no extra dependencies. Suppose
\texttt{orders.csv} has columns \texttt{order\_id,user\_id,amount}.

\begin{lstlisting}[style=pwshstyle]
# Build a lookup: user_id -> user object.
$userById = @{}
Import-Csv .\users.csv | ForEach-Object { $userById[$_.user_id] = $_ }

# Stream orders, attach the user's name and role to each order.
Import-Csv .\orders.csv |
    ForEach-Object {
        $u = $userById[$_.user_id]
        [PSCustomObject]@{
            order_id = [int]$_.order_id
            user_id  = [int]$_.user_id
            name     = $u.name
            role     = $u.role
            amount   = [double]$_.amount
        }
    } |
    Export-Csv -NoTypeInformation -Path .\orders_joined.csv -Encoding utf8
\end{lstlisting}

This is O(n+m) and streams orders one at a time -- so it handles
million-row order tables just fine as long as \texttt{users.csv} fits
in RAM. For a many-to-many join, the lookup value becomes a list and
you call \texttt{.Add()} on it.

\subsection*{Different delimiters and encodings}

\begin{lstlisting}[style=pwshstyle]
Import-Csv .\eu_users.csv -Delimiter ';'            # European-style CSV
Import-Csv .\data.tsv     -Delimiter "`t"           # tab-separated; backtick-t is the tab escape
Import-Csv .\latin1.csv   -Encoding latin1
\end{lstlisting}

\subsection*{When to drop into Python}

\cmd{Import-Csv} loads the whole file into memory as objects -- fine
for tens of thousands of rows, slow for millions. For anything in
\texttt{pandas} territory (joins, time-series resampling, vectorized
math), let PowerShell handle the orchestration and shell out to a
short Python script for the heavy work:

\begin{lstlisting}[style=pwshstyle]
# orchestrate from PowerShell, compute in pandas
python -m my_pipeline.transform --input .\users.csv --output .\users_enriched.csv
$enriched = Import-Csv .\users_enriched.csv
\end{lstlisting}

\begin{compare}
\cmd{Import-Csv} is to \texttt{pandas.read\_csv} what
\texttt{csv.DictReader} is to \texttt{pandas.read\_csv}: row-oriented,
object-per-row, all values as strings unless you cast. Use it for
small-to-medium files and orchestration; reach for pandas when you
need columnar speed.
\end{compare}

% =========================================================
\section{Data Wrangling -- JSON and HTTP}
\label{sec:json}
% =========================================================

\subsection*{Parsing and emitting JSON}

\begin{lstlisting}[style=pwshstyle]
# Parse a JSON file into a real object.
$config = Get-Content .\config.json -Raw | ConvertFrom-Json

# Or parse a JSON string you already have in a variable.
$obj = '{"id": 42, "tags": ["a", "b"]}' | ConvertFrom-Json
$obj.tags[0]            # 'a'

# Navigate with dot notation, just like in Python.
$config.database.host
$config.database.replicas[0].port

# Build an object and emit JSON (-Compress gives one line, no whitespace):
[PSCustomObject]@{ user = 'shiru'; roles = @('DE', 'MLE') } |
    ConvertTo-Json -Compress
# {"user":"shiru","roles":["DE","MLE"]}

# Modify and write back. -Depth defaults to 2, which is too shallow for
# almost any real config. Always pass -Depth explicitly when writing.
$config.database.host = 'mongo-prod-1'
$config | ConvertTo-Json -Depth 10 | Set-Content -Path .\config.json -Encoding utf8
\end{lstlisting}

\begin{gotcha}
\cmd{ConvertTo-Json}'s default \texttt{-Depth} is 2. Anything nested
deeper than that gets stringified as \texttt{"System.Collections.Hashtable"}
without warning. \textbf{Always pass \texttt{-Depth 10}} (or higher)
for any non-trivial structure.
\end{gotcha}

\subsection*{Hitting a REST API}

\cmd{Invoke-RestMethod} is one cmdlet that does
\texttt{requests.get(...).json()} in one step: it sends the HTTP
request \emph{and} parses the JSON response into a navigable object.

\begin{lstlisting}[style=pwshstyle]
# Simple GET.
$resp = Invoke-RestMethod -Uri 'https://api.github.com/repos/python/cpython'
$resp.stargazers_count
$resp.license.spdx_id

# GET with query parameters: pass a hashtable to -Body and PowerShell
# URL-encodes it onto the URL (?q=data+engineer&limit=10).
$resp = Invoke-RestMethod -Uri 'https://api.example.com/search' `
                          -Body @{ q = 'data engineer'; limit = 10 }

# GET with headers (e.g. bearer token from an env var).
$headers = @{ Authorization = "Bearer $env:GITHUB_TOKEN" }
$prs = Invoke-RestMethod -Uri 'https://api.github.com/repos/python/cpython/pulls?state=open' `
                         -Headers $headers
$prs | Select-Object number, title, user | Format-Table

# POST with a JSON body.
$body = @{ name = 'shiru'; role = 'DE' } | ConvertTo-Json -Depth 5
Invoke-RestMethod -Uri 'https://httpbin.org/post' `
                  -Method Post `
                  -ContentType 'application/json' `
                  -Body $body
\end{lstlisting}

The backtick at end-of-line is PowerShell's line-continuation
character (the equivalent of bash's trailing backslash). Use it
sparingly; most of the time you can just break inside parentheses
or at a pipe symbol.

\subsection*{Paginated APIs}

A practical pattern -- consume a paginated REST endpoint and stream
results into a single array.

\begin{lstlisting}[style=pwshstyle]
$all = @()
$page = 1
do {
    $items = Invoke-RestMethod `
        -Uri "https://api.example.com/items?page=$page&per_page=100" `
        -Headers @{ Authorization = "Bearer $env:API_TOKEN" }
    $all += $items
    $page++
} while ($items.Count -eq 100)

$all | ConvertTo-Json -Depth 10 |
    Set-Content -Path .\all_items.json -Encoding utf8
\end{lstlisting}

\begin{compare}
The PowerShell version is shorter than the Python equivalent because
\cmd{Invoke-RestMethod} bakes in JSON parsing and HTTP error handling.
The trade-off: no async, no session pooling, fewer auth strategies.
For high-throughput ingestion, use Python's \texttt{httpx} or
\texttt{aiohttp}.
\end{compare}

\subsection*{Error handling and retries}

\cmd{Invoke-RestMethod} throws a terminating error on non-2xx
responses. Wrap calls in \texttt{try / catch} and inspect
\texttt{\$\_.Exception.Response.StatusCode} to decide whether to
retry, skip, or fail loudly.

\begin{lstlisting}[style=pwshstyle]
function Invoke-WithRetry {
    param(
        [string]$Uri,
        [hashtable]$Headers = @{},
        [int]$MaxAttempts = 5
    )
    for ($attempt = 1; $attempt -le $MaxAttempts; $attempt++) {
        try {
            return Invoke-RestMethod -Uri $Uri -Headers $Headers
        } catch {
            $code = $_.Exception.Response.StatusCode.value__
            if ($code -in 429, 500, 502, 503, 504) {
                $backoff = [Math]::Pow(2, $attempt)
                Write-Warning "HTTP $code on attempt $attempt; sleeping $backoff s"
                Start-Sleep -Seconds $backoff
            } else {
                throw   # 4xx other than 429 -> don't retry
            }
        }
    }
    throw "Gave up after $MaxAttempts attempts: $Uri"
}

$data = Invoke-WithRetry -Uri 'https://api.example.com/items' `
                         -Headers @{ Authorization = "Bearer $env:API_TOKEN" }
\end{lstlisting}

This is the same exponential-backoff pattern you'd write around
\texttt{requests.get(...)} in Python, just in PowerShell's syntax.

% =========================================================
\section{Environment Variables and PATH}
\label{sec:env}
% =========================================================

This is where you wire up credentials and config for AWS, MongoDB,
MySQL, OpenAI, etc.

\subsection*{Reading and writing in the current session}

\begin{lstlisting}[style=pwshstyle]
$env:AWS_PROFILE                 # read (returns $null if unset)
$env:AWS_PROFILE = 'production'  # set for THIS process only
Remove-Item Env:\AWS_PROFILE     # unset for this process

# Inspect everything:
Get-ChildItem Env:
Get-ChildItem Env: | Where-Object Name -like 'AWS_*'

# Inspect PATH (it's just a string, semicolon-separated on Windows):
$env:Path -split ';'
\end{lstlisting}

\subsection*{Persisting across sessions}

\texttt{\$env:FOO = '...'} only affects the current process. To make a
variable stick across PowerShell restarts and reboots, write it to
the user or machine environment:

\begin{lstlisting}[style=pwshstyle]
# User-scoped (just you, all your shells): persists across sessions
[Environment]::SetEnvironmentVariable('AWS_PROFILE', 'production', 'User')

# Read it back from the persisted store (not from the current process):
[Environment]::GetEnvironmentVariable('AWS_PROFILE', 'User')

# Machine-scoped (every user on this PC): requires admin shell
[Environment]::SetEnvironmentVariable('CUDA_HOME', 'C:\CUDA\v12.5', 'Machine')

# Prepending to PATH (user-scoped):
$path = [Environment]::GetEnvironmentVariable('Path', 'User')
[Environment]::SetEnvironmentVariable('Path', 'C:\tools;' + $path, 'User')
\end{lstlisting}

\begin{gotcha}
Setting a User or Machine env var \emph{does not affect the current
PowerShell session.} You must open a new shell, or refresh
\texttt{\$env:Path} manually with
\texttt{\$env:Path = [Environment]::GetEnvironmentVariable('Path','Machine') + ';' + [Environment]::GetEnvironmentVariable('Path','User')}.
\end{gotcha}

\subsection*{A practical credentials pattern}

Don't paste credentials into scripts. Set them once at the user level,
read them in the script via \texttt{\$env:NAME}.

\begin{lstlisting}[style=pwshstyle]
# One-time setup at the user level:
[Environment]::SetEnvironmentVariable('AWS_ACCESS_KEY_ID',     'AKIA...', 'User')
[Environment]::SetEnvironmentVariable('AWS_SECRET_ACCESS_KEY', '...',     'User')
[Environment]::SetEnvironmentVariable('AWS_DEFAULT_REGION',    'us-east-1', 'User')
[Environment]::SetEnvironmentVariable('MONGO_URI',             'mongodb+srv://...', 'User')
[Environment]::SetEnvironmentVariable('MYSQL_HOST',            'db.internal', 'User')

# In every later script, just read:
$client = New-MongoClient -Uri $env:MONGO_URI
\end{lstlisting}

\begin{compare}
Unix shells use \bcmd{export FOO=bar} (process-scoped, with the
\texttt{.bashrc} / \texttt{.zshrc} dance for persistence). PowerShell
gives you both knobs explicitly: \texttt{\$env:FOO = ...} for the
process, \texttt{[Environment]::SetEnvironmentVariable(...)} for the
registry-backed persistent store.
\end{compare}

% =========================================================
\section{Python on Windows PowerShell}
% =========================================================

\subsection*{Where is Python?}

\cmd{Get-Command} is your \texttt{which}:

\begin{lstlisting}[style=pwshstyle]
Get-Command python
# CommandType  Name        Source
# Application  python.exe  C:\Users\shiru\miniconda3\python.exe

Get-Command python -All     # every python.exe on PATH, in order
\end{lstlisting}

If you have multiple Python installations (conda, the Microsoft Store
Python, the official installer\dots), the first one on
\texttt{\$env:Path} wins. \cmd{Get-Command python -All} shows them
all so you can reorder PATH or pin a specific path in a script.

\subsection*{The \texttt{py} launcher (optional)}

The official Python installer for Windows ships a \texttt{py.exe}
launcher that picks the right Python version from a single command.
It is \emph{not} installed by Conda or Microsoft Store, so it may not
be on your machine. If it is, you can do:

\begin{lstlisting}[style=pwshstyle]
py -0p             # list every Python interpreter the launcher knows about
py -3.12 script.py # run with Python 3.12 specifically
py -3 -m venv .venv
\end{lstlisting}

If you don't have \texttt{py}, just use \texttt{python} directly --
nothing in this tutorial requires the launcher.

\subsection*{Running scripts}

\begin{lstlisting}[style=pwshstyle]
python script.py                     # equivalent to bash: python3 script.py
python -m pytest                     # run a module as a script
python -m pip install requests       # always prefer `python -m pip` over bare `pip`
python -c "import sys; print(sys.executable)"
\end{lstlisting}

The \texttt{python -m} form is portable across all OSes and avoids any
ambiguity about which \texttt{pip} or \texttt{pytest} is on PATH.

% =========================================================
\section{Virtual Environments with \texttt{venv}}
% =========================================================

\subsection*{Create and activate}

\begin{lstlisting}[style=pwshstyle]
python -m venv .venv                 # create the venv in .\.venv
.\.venv\Scripts\Activate.ps1         # activate it (note: \Scripts\, not \bin\)
# Prompt becomes (.venv) PS C:\...>
python -m pip install requests pandas boto3 pymongo mysql-connector-python
deactivate                           # leaves the venv
\end{lstlisting}

\subsection*{The execution policy wall}

The first time you run \texttt{Activate.ps1} on a Windows machine,
you'll get this error:

\begin{lstlisting}[style=pwshstyle]
.\.venv\Scripts\Activate.ps1 : File C:\Users\shiru\Desktop\powershell\.venv\
Scripts\Activate.ps1 cannot be loaded because running scripts is disabled on
this system. ...
\end{lstlisting}

This is PowerShell's execution policy refusing to run unsigned local
scripts. Fix it once, at the user scope (\textbf{never machine
scope}), then forget about it:

\begin{lstlisting}[style=pwshstyle]
Set-ExecutionPolicy -Scope CurrentUser RemoteSigned
\end{lstlisting}

\texttt{RemoteSigned} means: local scripts run freely; scripts
\emph{downloaded from the internet} must be signed or unblocked. This
is the right balance for a developer machine, and it only affects
your user account.

\begin{compare}
There is no Unix equivalent. The execution policy is a Windows-only
trust mechanism. \texttt{Set-ExecutionPolicy -Scope CurrentUser
RemoteSigned} is the universally-recommended one-time fix.
\end{compare}

\subsection*{Worked example: per-project venv}

\begin{lstlisting}[style=pwshstyle]
# From a fresh project directory:
Set-Location C:\Users\shiru\projects\churn-model
python -m venv .venv
.\.venv\Scripts\Activate.ps1
python -m pip install --upgrade pip
python -m pip install -r requirements.txt
# work, work, work...
deactivate
\end{lstlisting}

To re-enter the venv tomorrow, the only command you need is
\texttt{.\textbackslash{}.venv\textbackslash{}Scripts\textbackslash{}Activate.ps1}.

% =========================================================
\section{\texttt{uv} on PowerShell}
% =========================================================

\texttt{uv} is the fast Python package and project manager you
already use. On Windows it installs without admin rights and
integrates with PowerShell exactly the same way as on Linux.

\subsection*{Install}

The official installer is a single PowerShell command:

\begin{lstlisting}[style=pwshstyle]
powershell -ExecutionPolicy ByPass -c "irm https://astral.sh/uv/install.ps1 | iex"
\end{lstlisting}

\texttt{irm} is the alias for \cmd{Invoke-RestMethod}; \texttt{iex}
is \cmd{Invoke-Expression}. Together they fetch the installer script
and run it. After install, open a new PowerShell tab so \texttt{uv}
appears on PATH, and verify:

\begin{lstlisting}[style=pwshstyle]
uv --version
\end{lstlisting}

\subsection*{Daily commands}

\begin{lstlisting}[style=pwshstyle]
# Start a new project (creates pyproject.toml, .venv, and uv.lock):
uv init churn-model
Set-Location churn-model

# Add dependencies (updates pyproject.toml and uv.lock automatically):
uv add pandas boto3 pymongo "mysql-connector-python>=8.3"
uv add --dev pytest ruff

# Run a Python script through uv (no manual activation needed):
uv run python -m churn_model.train --input data\users.csv

# Or run a one-off command in the project env:
uv run pytest
uv run ruff check .

# Sync the env to match uv.lock exactly (use this on CI / a fresh clone):
uv sync

# Lock without installing (useful in CI):
uv lock
\end{lstlisting}

\subsection*{When you do want to activate}

\texttt{uv run} usually means you never have to activate the venv.
But if you want a shell with the venv on PATH (e.g.\ for interactive
work), activate it like any other venv:

\begin{lstlisting}[style=pwshstyle]
.\.venv\Scripts\Activate.ps1
\end{lstlisting}

\subsection*{Working with \texttt{requirements.txt}}

If you have a legacy \texttt{requirements.txt}, \texttt{uv} has a fast
pip-compatible interface:

\begin{lstlisting}[style=pwshstyle]
uv venv                                       # create .venv
uv pip install -r requirements.txt           # fast pip install into .venv
uv pip compile requirements.in -o requirements.txt   # resolve & pin
uv pip sync requirements.txt                  # make .venv match requirements.txt EXACTLY
\end{lstlisting}

\begin{compare}
The mental model: \texttt{uv} is to \texttt{pip + venv + pip-tools +
virtualenv} what \texttt{cargo} is to Rust's loose tooling. \texttt{uv
run} replaces \texttt{source .venv/bin/activate \&\& python ...} as a
one-liner.
\end{compare}

% =========================================================
\section{Putting It Together: A Mini ETL Walkthrough}
% =========================================================

This section ties everything together. Goal: read a CSV of users, hit a
REST API to enrich each row with a profile lookup, run a small Python
transformation in a \texttt{uv}-managed venv to compute a per-user
score, write the result as JSON, log progress to a daily log file.

\subsection*{Project layout}

\begin{lstlisting}[style=pwshstyle,language=]
C:\projects\enrich\
    pyproject.toml         # managed by uv
    uv.lock
    .venv\                  # created by `uv venv`
    data\
        users.csv          # input
    enrich.ps1             # the orchestrator (PowerShell)
    score.py               # the transformation (Python)
    logs\
        enrich-2026-05-25.log
\end{lstlisting}

\subsection*{The Python transformation (\texttt{score.py})}

\begin{lstlisting}[style=pwshstyle,language=]
# score.py -- reads enriched rows from stdin (JSON lines),
# emits scored rows to stdout (JSON lines).
import json, sys

for line in sys.stdin:
    rec = json.loads(line)
    base = float(rec.get("score", 0))
    boost = 0.1 if rec.get("verified") else 0.0
    rec["final_score"] = round(base + boost, 4)
    sys.stdout.write(json.dumps(rec) + "\n")
\end{lstlisting}

\subsection*{The orchestrator (\texttt{enrich.ps1})}

\begin{lstlisting}[style=pwshstyle]
# enrich.ps1
param(
    [string]$Input  = '.\data\users.csv',
    [string]$Output = '.\data\users_scored.json'
)

$ErrorActionPreference = 'Stop'

# Set up a daily log file.
$logDir = '.\logs'
if (-not (Test-Path $logDir)) { New-Item -ItemType Directory $logDir | Out-Null }
$logFile = Join-Path $logDir ("enrich-{0:yyyy-MM-dd}.log" -f (Get-Date))

function Log {
    param([string]$Msg)
    $line = "{0:yyyy-MM-dd HH:mm:ss} {1}" -f (Get-Date), $Msg
    Add-Content -Path $logFile -Value $line
    Write-Output $line
}

Log "Reading $Input"
$users = Import-Csv $Input

Log "Enriching $($users.Count) users via API"
$enriched = foreach ($u in $users) {
    try {
        $profile = Invoke-RestMethod -Uri "https://api.example.com/users/$($u.user_id)" `
                                     -Headers @{ Authorization = "Bearer $env:API_TOKEN" }
        [PSCustomObject]@{
            user_id  = [int]$u.user_id
            name     = $u.name
            role     = $u.role
            score    = [double]$u.score
            verified = $profile.verified
        }
    } catch {
        Log "WARN: failed to enrich user_id=$($u.user_id): $_"
    }
}

Log "Scoring via Python in uv-managed venv"
$jsonLines = $enriched | ForEach-Object { $_ | ConvertTo-Json -Compress -Depth 10 }
$scored = $jsonLines | uv run python score.py | ForEach-Object {
    $_ | ConvertFrom-Json
}

Log "Writing $($scored.Count) records to $Output"
$scored | ConvertTo-Json -Depth 10 | Set-Content -Path $Output -Encoding utf8

Log "Done"
\end{lstlisting}

\subsection*{Run it}

\begin{lstlisting}[style=pwshstyle]
$env:API_TOKEN = 'sk-...'
.\enrich.ps1 -Input .\data\users.csv -Output .\data\users_scored.json
\end{lstlisting}

What this example shows, in one place:

\begin{itemize}[leftmargin=*]
\item Parameters with defaults (\texttt{param(...)}).
\item Error policy (\texttt{\$ErrorActionPreference = 'Stop'}) so an
unhandled error halts the script.
\item File-system setup with \cmd{Test-Path} + \cmd{New-Item}.
\item A small reusable function (\texttt{Log}).
\item CSV input via \cmd{Import-Csv}, with typed casting on the way in.
\item REST enrichment via \cmd{Invoke-RestMethod} with a bearer-token
header from \texttt{\$env}.
\item Hand-off to a Python script via \texttt{uv run python score.py},
streaming JSON lines through stdin/stdout.
\item JSON output with explicit \texttt{-Depth} and explicit
\texttt{-Encoding}.
\end{itemize}

% =========================================================
\section{Common Pitfalls}
% =========================================================

\textbf{Execution policy.} The first \texttt{Activate.ps1} run fails on
every fresh Windows install. Fix once with
\texttt{Set-ExecutionPolicy -Scope CurrentUser RemoteSigned}.

\textbf{Path quoting.} Use double quotes around paths with spaces.
The escape character is the backtick, not the backslash. To put a
literal \texttt{\$} inside a double-quoted string, prefix it with a
backtick: \texttt{"`\$5"}.

\textbf{\cmd{ConvertTo-Json} depth.} The default \texttt{-Depth} is 2.
Nested data silently turns into the string
\texttt{"System.Collections.Hashtable"}. Always pass
\texttt{-Depth 10} (or higher) for anything non-trivial.

\textbf{File encoding.} \cmd{Out-File} and \cmd{Set-Content} have
different defaults in PowerShell 5.1 vs.\ 7. Always pass
\texttt{-Encoding utf8} when the file will be read by non-Windows
tools.

\textbf{\texttt{Import-Csv} columns are strings.} Cast with
\texttt{[double]}, \texttt{[int]}, \texttt{[bool]},
\texttt{[datetime]} before any numeric or date comparison.

\textbf{\texttt{\$env:} doesn't persist.} Setting
\texttt{\$env:FOO = 'bar'} only affects the current process. Use
\texttt{[Environment]::SetEnvironmentVariable(...)} to persist.

\textbf{Array \texttt{+=} is O(n).} For large collections, use
\texttt{[System.Collections.Generic.List[object]]::new()} and call
\texttt{.Add()}, or just emit items through the pipeline and let it
do the right thing.

\textbf{Here-string closer at column 0.} The closing
\texttt{"@} or \texttt{'@} must be at the start of the line, with
no leading whitespace. Indenting it is a parse error.

\textbf{\texttt{-WhatIf} before destructive cmdlets.} Every
\texttt{Remove-}, \texttt{Stop-}, \texttt{Clear-} cmdlet accepts
\texttt{-WhatIf}. Make it a reflex.

\textbf{PowerShell 5.1 vs.\ 7.} On a modern Windows install, the
default \texttt{powershell.exe} is still PowerShell 5.1; the
cross-platform \texttt{pwsh.exe} is PowerShell 7+. Anything in this
tutorial that uses \texttt{utf8NoBOM} or \texttt{??} or the ternary
operator is PowerShell 7 only. Always launch \texttt{pwsh}, not
\texttt{powershell}, for new work.

% =========================================================
\section{Quick Reference: Bash $\to$ PowerShell}
% =========================================================

The minimal table you need pinned somewhere. PowerShell aliases
(\texttt{ls}, \texttt{cd}, \texttt{cat}, \dots) work in the interactive
shell but use the cmdlet name in scripts.

\begin{center}
\renewcommand{\arraystretch}{1.15}
\begin{tabular}{p{0.42\linewidth} p{0.50\linewidth}}
\toprule
\textbf{Bash} & \textbf{PowerShell} \\
\midrule
\bcmd{pwd}                              & \cmd{Get-Location} \\
\bcmd{cd /path}                          & \cmd{Set-Location} \texttt{C:\textbackslash{}path} \\
\bcmd{ls}                                & \cmd{Get-ChildItem} \\
\bcmd{ls -la}                            & \cmd{Get-ChildItem} \texttt{-Force} \\
\bcmd{find . -name "*.csv"}              & \cmd{Get-ChildItem} \texttt{-Recurse -Filter *.csv} \\
\bcmd{cat file}                          & \cmd{Get-Content} \texttt{file} \\
\bcmd{head -n 10 file}                   & \cmd{Get-Content} \texttt{file -TotalCount 10} \\
\bcmd{tail -n 10 file}                   & \cmd{Get-Content} \texttt{file -Tail 10} \\
\bcmd{tail -f file}                      & \cmd{Get-Content} \texttt{file -Wait -Tail 0} \\
\bcmd{wc -l file}                        & \texttt{(\cmd{Get-Content} file | \cmd{Measure-Object} -Line).Lines} \\
\bcmd{grep PATTERN file}                 & \cmd{Select-String} \texttt{-Pattern PATTERN -Path file} \\
\bcmd{which python}                      & \texttt{(\cmd{Get-Command} python).Source} \\
\bcmd{mkdir dir}                         & \cmd{New-Item} \texttt{-ItemType Directory dir} \\
\bcmd{mkdir -p a/b/c}                    & \cmd{New-Item} \texttt{-ItemType Directory -Force a\textbackslash{}b\textbackslash{}c} \\
\bcmd{touch f}                           & \texttt{if (-not (Test-Path f)) \{ \cmd{New-Item} -ItemType File f \}} \\
\bcmd{cp src dst}                        & \cmd{Copy-Item} \texttt{src dst} \\
\bcmd{cp -r src dst}                     & \cmd{Copy-Item} \texttt{-Recurse src dst} \\
\bcmd{mv src dst}                        & \cmd{Move-Item} \texttt{src dst} \\
\bcmd{rm f}                              & \cmd{Remove-Item} \texttt{f} \\
\bcmd{rm -rf dir}                        & \cmd{Remove-Item} \texttt{-Recurse -Force dir} \\
\bcmd{ln -s tgt link}                    & \cmd{New-Item} \texttt{-ItemType SymbolicLink -Path link -Target tgt} \\
\bcmd{echo \$VAR}                        & \texttt{\$env:VAR} \\
\bcmd{export VAR=x}                      & \texttt{\$env:VAR = 'x'} \quad (process-scoped) \\
\bcmd{export VAR=x} (persistent)          & \texttt{[Environment]::SetEnvironmentVariable('VAR','x','User')} \\
\bcmd{cmd > out}                          & \texttt{cmd | \cmd{Set-Content} out} \\
\bcmd{cmd >> out}                         & \texttt{cmd | \cmd{Add-Content} out} \\
\bcmd{cmd1 \&\& cmd2}                     & \texttt{cmd1 \&\& cmd2} \quad (PowerShell 7+) \\
\bcmd{cmd1 || cmd2}                       & \texttt{cmd1 || cmd2} \quad (PowerShell 7+) \\
\bcmd{a=\$(cmd)}                          & \texttt{\$a = (cmd)} \\
\bcmd{curl URL | jq .}                    & \texttt{\cmd{Invoke-RestMethod} URL} \\
\bcmd{2>/dev/null}                        & \texttt{2>\$null} \\
\bcmd{for f in *.csv; do \dots done}      & \texttt{\cmd{Get-ChildItem} *.csv | \cmd{ForEach-Object} \{ \dots \}} \\
\bottomrule
\end{tabular}
\end{center}

\vspace{1em}
\noindent\textit{One last thing:} when you're stuck, start with
\cmd{Get-Command} \texttt{*word*} (find the cmdlet), \cmd{Get-Help}
\texttt{NAME -Examples} (see how to use it), and \cmd{Get-Member}
(see what properties an object has). These three make PowerShell easy
to explore on your own, without leaving the shell.

\end{document}
